\section{Introduction}\label{sec:introduction}
QLang is a high-level programming language designed for quantum computing and distributed quantum networking. 
It provides a simple and intuitive syntax for expressing quantum algorithms, operations, and communication between separate quantum nodes called places, 
making it accessible to both beginners and experienced programmers in the field of quantum computing.

% ─────────────────────────────────────────────

\subsection{Overview}\label{sec:overview}
A QLang program is made up of top-level items: \texttt{place} declarations,
function declarations, import directives, and statements. A place is QLang's
primary unit of concurrent execution. Each place has its own scope, defines
its own functions, and runs as a separate process. Code written outside any
named place belongs to the implicit \texttt{global} place, 
which only executes if it contains at least one statement
or function declaration.

The language supports five built-in variable types: \texttt{state},
\texttt{obs}, \texttt{num}, \texttt{text}, and \texttt{any}. Quantum-specific operations,
including the \texttt{measure} keyword and all built-in gates, are first-class constructs 
in QLang, not library calls. The language natively supports quantum gates such as 
\texttt{superpose}, \texttt{shift}, \texttt{not}, \texttt{dual\_not}, \texttt{phase\_not}, 
\texttt{entangle}, \texttt{entangle\_phase}, and \texttt{swap}.
QLang is interpreted. Source files with the extension \texttt{.ql}
are parsed by an ANTLR4-generated parser and executed by the QLang interpreter
written in Python.

% ─────────────────────────────────────────────

\subsection{Quick Example}\label{sec:quick-example}
The following complete QLang program creates a Bell state (a maximally
entangled two-qubit state), measures both qubits, and prints the results:
% Auto-generated by docs/tools/highlight.py - do not edit
\begin{codebox}
(*@\textcolor{codebox_comment}{place}@*) (*@\textcolor{black}{Alice}@*) (*@\textcolor{codebox_operator}{\{}@*)
    (*@\textcolor{codebox_type}{\textbf{state}}@*) (*@\textcolor{black}{a}@*)(*@\textcolor{codebox_operator}{,}@*) (*@\textcolor{black}{b}@*)(*@\textcolor{codebox_operator}{;}@*)

    (*@\textcolor{codebox_builtin_operator}{H}@*) (*@\textcolor{black}{a}@*)(*@\textcolor{codebox_operator}{;}@*)           (*@\textcolor{codebox_comment}{// Put qubit \string'a\string' into superposition}@*)
    (*@\textcolor{codebox_builtin_operator}{CNOT}@*) (*@\textcolor{black}{a}@*) (*@\textcolor{codebox_operator}{->}@*) (*@\textcolor{black}{b}@*)(*@\textcolor{codebox_operator}{;}@*)   (*@\textcolor{codebox_comment}{// Entangle \string'a\string' with \string'b\string' — this creates a Bell state}@*)

    (*@\textcolor{codebox_type}{\textbf{obs}}@*) (*@\textcolor{black}{ma}@*) (*@\textcolor{codebox_operator}{=}@*) (*@\textcolor{codebox_builtin_operator}{measure}@*) (*@\textcolor{black}{a}@*)(*@\textcolor{codebox_operator}{;}@*)
    (*@\textcolor{codebox_type}{\textbf{obs}}@*) (*@\textcolor{black}{mb}@*) (*@\textcolor{codebox_operator}{=}@*) (*@\textcolor{codebox_builtin_operator}{measure}@*) (*@\textcolor{black}{b}@*)(*@\textcolor{codebox_operator}{;}@*)

    (*@\textcolor{codebox_number}{println}@*)(*@\textcolor{codebox_operator}{(}@*)(*@\textcolor{codebox_string}{\string"Measurement A: \%d\string"}@*) (*@\textcolor{codebox_operator}{\%}@*) (*@\textcolor{black}{ma}@*)(*@\textcolor{codebox_operator}{)}@*)(*@\textcolor{codebox_operator}{;}@*)
    (*@\textcolor{codebox_number}{println}@*)(*@\textcolor{codebox_operator}{(}@*)(*@\textcolor{codebox_string}{\string"Measurement B: \%d\string"}@*) (*@\textcolor{codebox_operator}{\%}@*) (*@\textcolor{black}{mb}@*)(*@\textcolor{codebox_operator}{)}@*)(*@\textcolor{codebox_operator}{;}@*)
    (*@\textcolor{codebox_comment}{// Both will always agree: both 0 or both 1}@*)
(*@\textcolor{codebox_operator}{\}}@*)

\end{codebox}

After running, both measurements will always agree, both \texttt{0} or both
\texttt{1}, which is the defining property of a Bell state.

% ─────────────────────────────────────────────

\subsection{Design Principles}\label{sec:design-principles}
\subsubsection{Classical and Quantum Code in One Language}
QLang does not separate quantum and classical concerns into different APIs or
files. A variable of type \texttt{obs} and a variable of
type \texttt{state} (qubit) can appear side-by-side in the same scope.
Measurement, when invoked on a \texttt{state}, collapses the qubit and returns an 
\texttt{obs} value indicating the outcome: \texttt{F} (false) if the qubit collapsed to |0⟩, 
and \texttt{T} (true) if it collapsed to |1⟩.
\subsubsection{Structured Error Reporting} 
When the interpreter encounters an error, whether a syntax mistake, an
undefined variable, or an illegal runtime operation, it displays a detailed
error box that highlights the exact source location, shows surrounding context
lines, and explains the cause.
\subsubsection{Modular Structure via Places}
QLang organizes code into units called \texttt{places}. Each \texttt{place} has its own 
name, its own variables, and runs independently from other places. Places can communicate 
with each other using the built-in \texttt{send} and \texttt{receive} instructions, which 
transfer both classical and quantum data. This makes it natural to write distributed 
programs, such as quantum communication protocols, directly in QLang.